\section{结论}
本文针对规避者需同时规避两枚非合作飞行器并保持对指定目标终端精度的需求，提出了一种三维一体化规避制导框架。建立了四体相对运动模型，用以描述规避者、目标与两枚非合作飞行器之间的耦合交互。在此基础上，将双非合作飞行器规避需求、目标到达精度与控制能耗惩罚纳入统一性能指标，推导了最优闭环制导律。该制导律显式计入交叉耦合效应，即一次机动同时影响两次规避交会与终端目标到达。

为在不同交会几何下自适应选择制导权重，提出了基于BC-TD3的参数优化策略。利用单打靶优化方法生成的专家示范对策略进行初始化，再经TD3微调提升其在非线性多约束环境中的适应能力。此外，引入带稀疏正则化与遗忘因子的递推最小二乘辨识器，从含噪测量中估计状态空间模型的非平凡参数，从而降低制导回路对测量误差的敏感性。

仿真结果表明，所提方法能够在三维双非合作飞行器场景中可靠地同时完成规避与目标到达。在无噪声蒙特卡洛试验中，BC-TD3的规避成功率与目标到达成功率分别为94.7\%和94.3\%，优于传统比例导引、程序正弦机动以及不含行为克隆的TD3。在测量噪声条件下，引入稀疏在线辨识后，两项成功率进一步提高至96.7\%和96.0\%。上述结果表明，所提框架可为多非合作飞行器一体化规避与终端制导提供有效且鲁棒的解决方案。后续工作将进一步考虑更贴近实际的飞行器动力学、异构及机动非合作飞行器、执行机构约束，以及更广泛不确定性条件下的验证。
